%==============================================================
%  General Conclusion
%==============================================================

When this project began, the central question was whether a single student, working alone over four months, could build \deepguard{}, a complete platform capable of detecting deepfakes in images and videos and presenting the results in a way that non-experts could actually use. The work documented in this report shows that the answer is yes. The platform exists, the model performs accurately on its target inputs, the forensic analysis runs end to end, and the user interface meets the standards of a modern web application. The problem that motivated the work --- the erosion of trust in visual media --- has not disappeared, but the kind of accessible tool needed to push back against it is now one step closer to existence.

The technical scope covered every layer of a production-style web application. The AI side combines a Vision Transformer model from Hugging Face with MTCNN face detection, an image normalization pipeline, and a risk scoring system that translates raw probabilities into a four-level categorical verdict. The backend, built on FastAPI, exposes a REST API protected by JWT authentication where appropriate, and persists every result in a PostgreSQL database through SQLAlchemy. The frontend is a React single-page application with eight pages, three theme modes, drag-and-drop file uploads, animated charts, and CSV export. The forensic module reads EXIF metadata and flags suspicious patterns, and the PDF report generator produces a downloadable forensic document for every scan. Every feature listed in the original specification was delivered.

The seven objectives stated in Chapter \ref{chap:context} of this report can now be reviewed against what was actually built. The deepfake detection engine handles both images and videos and reached the accuracy thresholds set out in the initial validation. The risk scoring system translates the model output into the four-level verdict expected by the user. The forensic module extracts EXIF metadata and flags missing or suspicious fields. The backend authenticates users through JWT and isolates data per account. The frontend supports drag-and-drop uploads, dark mode, and the responsive design promised in the requirements. The PDF reports and shareable links work as specified. And the codebase is organized into single-responsibility modules that follow a clean, documented structure. All seven objectives were met.

The development was not free of difficulties. The most instructive challenge came from the image normalization step. Early versions of the detector were too sensitive to photo filters and social-media-style color adjustments, which pushed genuine photographs into the suspicious range and undermined the credibility of the verdicts. Discovering this and building a dedicated normalization pipeline took longer than expected and reshaped the way the detection logic was assembled. Other challenges were smaller but accumulated over time: the choice between Scrum and 2TUP at the start of the project, the decision to sample every tenth frame instead of every frame in videos to keep the pipeline fast, the cover page of this very report which required several iterations before fitting on a single page, and the constant calibration between scope and quality that any solo project demands. None of these were blockers, but each one influenced the final shape of the platform in ways worth recording.

Several directions could extend the work presented here. The most important is retraining the dima806 model on a more recent labeled dataset. The current model was trained roughly three years ago, and the generative AI tools that produce deepfakes today create images with visual artifacts very different from those the model originally learned to recognize. This concept drift is the single largest threat to the platform's long-term accuracy, and the model's own author has explicitly recommended periodic retraining as a mitigation. A second direction concerns the scope of the analysis itself. The current model is specialized for faces, and adding a second model trained on full-body figures, manipulated landscapes, or AI-generated documents would broaden the platform's coverage substantially without disrupting the existing pipeline. Beyond these two priorities, lighter improvements are also possible: a real-time analysis mode for live video streams, a multi-language user interface (starting with French and Arabic for the regional context), a native mobile application that reuses the same backend API, and production-grade features such as rate limiting and observability for a public deployment.

Beyond its technical results, the project offered a chance to build a complete software product from scratch, to make engineering decisions and live with their consequences, and to write about all of it in a way that explains rather than impresses. The final platform is by no means a definitive answer to the deepfake problem, but it is a credible, accessible, and honest contribution to a debate that is only going to grow louder in the years ahead.